% Generated draft: requires mathematical and source review.
\documentclass[11pt]{article}
\usepackage[margin=25mm]{geometry}
\usepackage{fontspec}
\setmainfont{Latin Modern Roman}
\usepackage{amsmath,amssymb,mathtools,unicode-math}
\setmathfont{Latin Modern Math}
\usepackage{enumitem,needspace,xurl,hyperref}
\setlength{\parindent}{0pt}
\setlength{\parskip}{5pt}
\setlength{\emergencystretch}{3em}
\newcommand{\N}{\mathbb{N}}
\newcommand{\Z}{\mathbb{Z}}
\newcommand{\Q}{\mathbb{Q}}
\newcommand{\R}{\mathbb{R}}
\newcommand{\C}{\mathbb{C}}
\newcommand{\F}{\mathbb{F}}
\newcommand{\A}{\mathbb{A}}
\newcommand{\PP}{\mathbb{P}}
\newcommand{\E}{\mathbb{E}}
\newcommand{\eps}{\varepsilon}
\newcommand{\dd}{\,\mathrm{d}}
\begin{document}
\section*{UnsolvedMath ID: 30006337}
\textbf{Source problem code:} catalogue entry
\subsection{Definitions and mathematical statement}
Let \(q\) be an indeterminate and put \(A=\Z[[q-1]]\). For a smooth \(\Z\)-algebra \(R\), write \(\Omega^\bullet_{R/\Z}\) for its de Rham complex, with Hodge filtration
\[
F^i\Omega^\bullet_{R/\Z}=\Omega^{\geq i}_{R/\Z}.
\]
The induced filtration on de Rham cohomology is
\[
F^iH^n_{\mathrm{dR}}(R/\Z)
=\operatorname{im}\!\left(
H^n(F^i\Omega^\bullet_{R/\Z})\longrightarrow
H^n(\Omega^\bullet_{R/\Z})\right).
\]
The supplied questions do not specify their full class of geometric inputs. Smooth \(\Z\)-algebras are explicitly mentioned in the accompanying qualification, so the notation here describes that setting without asserting that it is the complete intended scope.

\textbf{Question A: deformation of the Hodge filtration.}
Does the Hodge filtration admit a \(q\)-deformation on the intended \(q\)-de Rham object? At the complex level, let \(K_R=q\text{-}\Omega_R\) denote that object over \(A\). A filtered deformation would require a decreasing filtration \(F_q^\bullet K_R\) and a filtered specialization equivalence
\[
\left(F_q^\bullet K_R\right)\otimes_A^{\mathbf L} A/(q-1)
\simeq
\left(F^\bullet\Omega^\bullet_{R/\Z}\right).
\]
This equation gives a complex-level interpretation of the question; the supplied record does not define \(K_R\), specify the relevant category or completion conventions, or decide whether the requested filtration is on the complex or only on its cohomology. Any additional requirements of canonicity, functoriality, multiplicativity, or completeness must therefore be specified before this becomes a fully determined existence problem.

\textbf{Question B: Habiro descent.}
The supplied coefficient ring is written
\[
\mathcal H=\varprojlim_{m\in\N}
\Z[q]^{\wedge}_{(q^m-1)}.
\]
For a positive integer \(m\), the displayed completion means
\[
\Z[q]^{\wedge}_{(q^m-1)}
=\varprojlim_{k\geq1}\Z[q]/\bigl((q^m-1)^k\bigr).
\]
The record does not specify the indexing order or transition maps for the outer inverse limit, or its convention for \(\N\); its formula is retained subject to that qualification. Once this inverse system and the intended coefficient map \(\iota:\mathcal H\to A\) are fixed, the question asks whether \(q\)-de Rham cohomology descends along \(\iota\). A complex-level interpretation asks for an object \(D_R\) over \(\mathcal H\), together with an equivalence
\[
D_R\otimes^{\mathbf L}_{\mathcal H,\iota}A\simeq K_R,
\]
or the corresponding completed derived base-change equivalence if the intended category requires completion. The source does not specify which interpretation is intended, whether descent concerns individual cohomology groups or a complex, or which structures and naturality must descend.

\textbf{Supplied qualifications.}
The original excerpt states that Question A has a negative answer in general. It also announces, as upcoming work, that for every smooth \(\Z\)-algebra \(R\) such that every prime
\[
p\leq\dim(R/\Z)
\]
is invertible in \(R\), the complex \(q\text{-}\Omega_R\) admits a canonical \(q\)-deformation of the Hodge filtration. The supplied dated literature assessment further reports that canonical \(q\)-Hodge filtrations exist after inverting finitely many small primes and that Habiro descent is proved whenever such a filtration exists. These are source-reported qualifications, not independently verified assertions here. The record labels the problem partially solved and says that the unrestricted questions are not fully settled; it does not identify a precise remaining conjecture or all hypotheses of the reported descent result.
\subsection{Short English statement}
The record asks whether the Hodge filtration has a \(q\)-deformation and whether \(q\)-de Rham cohomology descends from \(\Z[[q-1]]\) to the displayed Habiro ring. It reports a general obstruction to the first question, a canonical construction under a small-prime invertibility hypothesis, and conditional descent. The exact remaining target, descent category, and inverse-limit conventions require clarification.
\subsection{Sources}
\begin{itemize}
\item Supplied local UnsolvedMath registry record, permanent ID \(30006337\), title \emph{q-Hodge Filtrations and Habiro Descent}, exact source problem code \texttt{catalogue entry}. The record belongs to the Oberwolfach Reports Open Problems collection and gives proposed year 2025; no individual proposer is supplied.
\item \emph{Non-Archimedean Geometry and Applications} (2025), Oberwolfach Reports, pp.~1517--1578, DOI \texttt{10.4171/owr/2025/29}: \url{https://doi.org/10.4171/owr/2025/29}. This is the supplied report attribution and source link. No problem-specific page, speaker, or author attribution is provided.
\item Supplied literature source: \emph{Derived q-Hodge complexes and refined TC}, \url{https://arxiv.org/abs/2410.23115}.
\item Supplied literature source: \emph{Prisms and prismatic cohomology}, \url{https://arxiv.org/abs/1905.08229}.
\item Literature assessment dated exactly \texttt{2026-08-22}: the supplied curation labels the problem partially solved, reports failure of a general \(q\)-deformation, existence of canonical filtrations after suitable small primes are inverted, and Habiro descent whenever such a filtration exists. It says the unrestricted questions are not fully settled. This is supplied dated curation derived from the OWR record, not an independently verified literature assessment; the cited claims and their current status require independent verification.
\end{itemize}
\end{document}
